%% ---------------------------------------------------------------------
%  CHƯƠNG 6 — KẾT LUẬN VÀ HƯỚNG PHÁT TRIỂN
% ---------------------------------------------------------------------
\chapter{Kết luận và hướng phát triển}
\label{ch:ketluan}

\begin{muctieu}
Chương này tổng kết hệ thống đã xây dựng dựa trên source VisionAI Pro,
rút ra bài học và đề xuất các hướng phát triển phù hợp với kiến trúc hiện tại.
\end{muctieu}

\section{Kết luận}

Đồ án đã xây dựng một ứng dụng trợ lý thị giác AI bằng Python và Streamlit,
cho phép tải/chụp ảnh, phân tích bằng Gemini Vision và đọc văn bản cục bộ bằng
Tesseract. Ứng dụng bổ sung hỏi đáp theo ảnh, tìm vật, hỗ trợ di chuyển từ ảnh
tĩnh, lệnh thoại và phản hồi bằng giọng nói. SpeechRecognition/Google Web Speech
được dùng cho STT; Edge TTS và Windows SAPI5 đảm nhiệm TTS. Robot Vision, chế độ
trợ năng, lịch sử có lọc/yêu thích, thống kê, xuất TXT và kiểm tra kết nối
Gemini hoàn thiện các luồng sử dụng chính.

Kiến trúc module tách Vision, OCR, Speech, trợ năng, lệnh thoại và robot.
Trên source bàn giao, ứng dụng khởi chạy, health check đạt và 15/15 kiểm thử
cục bộ đạt. Đây là kết quả kiểm thử phần mềm, không phải số đo độ chính xác
của AI; Gemini API thật và chất lượng sử dụng với người khiếm thị chưa được
đánh giá. Cá nhân hoá hiện chỉ ở mức cấu hình và trạng thái phiên làm việc.

\section{Bài học rút ra}

Quá trình xây dựng cho thấy một ứng dụng AI hỗ trợ người khiếm thị không chỉ
cần mô hình AI mà còn cần thiết kế luồng dữ liệu, xử lý lỗi, khả năng tiếp cận,
phản hồi âm thanh và chỉ dẫn an toàn. Việc tách OCR cục bộ, nhận dạng giọng nói
trực tuyến, TTS, trợ năng và giao diện robot thành các module riêng giúp dễ
kiểm thử và thay thế công nghệ.

Đồ án cũng củng cố kỹ năng Python, Streamlit, xử lý ảnh, tích hợp API AI,
xử lý âm thanh, quản lý biến môi trường và thiết kế một pipeline đa phương thức.

\section{Hướng phát triển}

Một số hướng phát triển phù hợp với source hiện tại gồm:
\begin{itemize}[leftmargin=1.6em, itemsep=2pt]
  \item Xây dựng hồ sơ người dùng và cơ chế cá nhân hoá lâu dài thay cho cấu hình chỉ tồn tại trong phiên.
  \item Bổ sung bộ dữ liệu kiểm thử có nhãn để đo định lượng OCR, Vision, STT và khả năng nhận diện tiền.
  \item Nghiên cứu STT cục bộ hoặc các lựa chọn nhận dạng khác để giảm phụ thuộc vào dịch vụ mạng, đồng thời đánh giá quyền riêng tư.
  \item Mở rộng khả năng xử lý ảnh/video thời gian thực cho cảnh báo vật cản.
  \item Bổ sung cơ chế cache, theo dõi lỗi mạng và phương án vận hành phù hợp khi Gemini, Google Web Speech hoặc TTS trực tuyến không khả dụng.
  \item Đánh giá khả năng tiếp cận, lệnh thoại, kích thước chữ/nút và nội dung cảnh báo với người dùng mục tiêu.
  \item Phát triển phiên bản mobile hoặc thiết bị đeo sau khi hoàn thiện pipeline và kiểm thử với người dùng thực tế.
\end{itemize}
